\section{Linear operators and observables (9/2)}

\referto{\nielsen~2.1; \tong~3.2, 4.1; \sakurai~1.3-1.4; \shankar~1.5-1.6}

Given two vector spaces $\mathbb{V}$ and $\mathbb{W}$, we can also define a linear operator $\hat O$ as any function that maps the vector space $\mathbb{V}$ to $\mathbb{W}$, which is linear in its inputs:
\begin{align}
\hat O(\sum_ia_i\dket{v_i})=\sum_ia_i\hat O(\dket{v_i})
\end{align}
In the case of $\mathbb{V}=\mathbb{W}$ we call it a linear operator $\hat O\in\mathbb{L}_\mathbb{V}$ on vector space $\mathbb{V}$, which is of particular interest to quantum mechanics.

Using a complete set of orthonormal basis $\{\dket{j}|1\leq j\leq\text{dim}\mathbb{V}\}$ for Hilbert space $\mathbb{V}$, each linear operator has a matrix representation:
\begin{align}
\hat O=\sum_{a,b}O_{a,b}\dket{a}\dbra{b},~~~O_{a,b}\equiv\expval{a|\hat O|b}
\end{align}
In particular, for the identity matrix which maps every vector to itself, we have the \textbf{completeness relation} for basis $\{\dket{i}\}$, also known as the \textbf{resolution of identity}:
\begin{align}
\sum_i\dket{i}\dbra{i}=\hat 1.
\end{align}
In the case of continuous variables, e.g. coordinates for a particle in free space, the summation will be replaced by integral:
\begin{align}
    \int\dif x\dket{x}\dbra{x}=\hat1
\end{align}
where $\dket{x}$ is an eigenstate of position operator $\dket{x}$ with eigenvalue $x$. 

We can also define the product of two linear operators $\hat M,\hat N\in\mathbb{L}_\mathbb{V}$:
\begin{align}
\hat M\hat N=\sum_{a,b,c}M_{a,b}N_{b,c}\dket{a}\dbra{c}
\end{align}
their commutator
\begin{align}
[\hat M,\hat N]=\hat M\hat N-\hat N\hat M
\end{align}
and anti-commutator
\begin{align}
\{\hat M,\hat N\}=\hat M\hat N+\hat N\hat M
\end{align}
and the \emph{Hermitian conjugate} (or \emph{Hermitian adjoint}) of an operator
\begin{align}
\hat O^\dagger=\sum_{a,b}O^\ast_{b,a}\dket{a}\dbra{b}=(O^T)^\ast=(O^\ast)^T
\end{align}
Specifically in Postulate \ref{postulate 3:projective measurements} about quantum (projective) measurements, each physical observable corresponds to a Hermitian linear operator, represented as a Hermitian matrix in the orthonormal basis:
\begin{align}
\hat M=\hat M^\dagger=\sum_{a,b}M_{a,b}\dket{a}\dbra{b},~~~(M^\dagger)_{a,b}=M^\ast_{b,a}=M_{a,b}\equiv\expval{a|\hat M|b}.
\end{align}

The diagonal form
\begin{align}\label{diagonal form}
\hat O=\sum_a O_a\dket{a}\dbra{a}
\end{align}
where $\{\dket{a}\}$ is a complete set of orthonormal basis of the Hilbert space, is also known as the \emph{spectral decomposition} of a matrix $\hat O$. If there is a complete set of orthonormal basis $\{\dket{a}\}$ such that the linear operator $\hat O$ can be rewritten as the diagonal form (\ref{diagonal form}), the matrix $\hat O$ is called ``diagonalizable''.

Not all matrices are diagonalizable (i.e. have a diagonal form), two non-diagonalizable examples are
\begin{align}
\hat M_1=\bpm1&1\\0&1\epm,~\hat M_2=\bpm1&1\\0&0\epm,~~~
\end{align}
$\hat M_1$ is known as a Jordan block with only 1 eigenvector, while $\hat M_2$ has two eigenvectors with different eigenvalues, but not orthogonal to each other. Generally, any complex square matrix can be decomposed into tensor sums of Jordan blocks by similarity transformations.\\

Below are a few useful theorems on the diagonalization of matrices:

\begin{thm}
    \textbf{The spectral decomposition theorem}: A linear operator $\hat O$ on vector space $\mathbb{V}$ is diagonal w.r.t. some orthonormal basis of $\mathbb{V}$ if and only if operator $\hat O$ is \emph{normal} i.e. $\hat O^\dagger\hat O=\hat O\hat O^\dagger$. 
\end{thm}
Refer to \nielsen~Box 2.2 for proof of this theorem.\\

\begin{thm}
    A normal matrix is Hermitian if and only if it has real eigenvalues.
\end{thm}

\begin{thm}
    The eigenkets corresponding to different eigenvalues of a Hermitian matrix are orthogonal.
\end{thm}

\begin{thm}
   A positive linear operator must be Hermitian. 
\end{thm}

\begin{thm}
    \textbf{Simultaneous diagonalization theorem}: Two Hermitian linear operators $\hat A$ and $\hat B$ can be diagnoalized simultaneously if and only if $[\hat A,\hat B]=0$.
\end{thm}
 In other words, there is a complete set of simultaneous eigenket $\{\dket{a,b}\}$ such that
\begin{align}
\hat A\dket{a,b}=a\dket{a,b},~~~\hat B\dket{a,b}=b\dket{a,b}.
\end{align}
if and only if $[\hat A,\hat B]=0$.\\

An Hermitian operator can be diagonalized into a special orthonormal basis $\{\dket{i}\}$, known as its eigenbasis, in which the linear operator is represented by a diagonal matrix
\begin{align}
\label{eigen basis}
\hat M=\sum_i\lambda_i\dket{i}\dbra{i}
\end{align}
If we use the set $\{m\}$ to label all different (real) eigenvalues of matrix $\hat M$, we can rewrite the above diagonal form (\ref{eigen basis}) as
\begin{align}
\hat M=\sum_m m\hat P_m,~~~\hat P_m=\sum_{j}\delta_{\lambda_j,m}\dket{j}\dbra{j}
\end{align}
where operators $\hat P_m$ are known as \emph{projection operators} (or \emph{projectors}) onto the eigenspace of $\hat M$ with eigenvalue $m$. Mathematically, a projector $\hat P$ is defined as a linear operator satisfying:
\begin{align}
\hat P^2=\hat P
\end{align}
You will learn more about projectors in HW2.

Given a state $\dket{\psi}$ and a physical observable $\hat M$, the (probability-weighted) \emph{expectation value} $\expval{\hat M}$ of operator $\hat M$ w.r.t. state vector $\dket{\psi}$ is
\begin{align}
\label{expval}
\expval{\hat M}\equiv\frac{\expval{\psi|\hat M|\psi}}{\expval{\psi|\psi}}
\end{align}
Naturally, the expectation value of any physical observable ia a real number, since any physical observable corresponds to a Hermitian operator.

We note that the expectation value (\ref{expval}) of any observable is the same for two linearly dependent states $\dket{\psi}$ and $\alpha\dket{\psi}$ for any complex scalar $\alpha\in\mathbb{C}$. This means $\dket{\psi}$ and $\alpha\dket{\psi}$ describe the same physical state. In other words, in a $N$-dimensional Hilbert space $\mathbb{V}=\mathbb{C}^N$, for a generic state
\begin{align}
\dket\psi\equiv\sum_{j=1}^N\psi_j\dket{j}
\end{align}
where $\{\dket{j}|1\leq j\leq N\}$ is a complete set of orthonormal basis of $\mathbb{V}$, although $\{\psi_j\in\mathbb{C}|1\leq j\leq N\}$ correspond to $2N$ independent real variables, only $2N-2$ of them are physically measurable, because the global phase factor and normalization leads to 2 redundancies that cannot be physically observed.

In the simplest example of a single qubit, we have $\mathbb{V}=\mathbb{C}^2$ for the $N=2$-dimensional Hilbert space. There are 4 linearly independent Hermitian operators in $\mathbb{C}^2$, whose matrix representations are
\begin{align}
\hat 1=\bpm1&0\\0&1\epm,~~~\hat\sigma_x=\bpm0&1\\1&0\epm,~~~\hat\sigma_y=\bpm0&-\imth\\ \imth&0\epm,~~~\hat\sigma_z=\bpm1&0\\0&-1\epm.
\end{align}
As you will learn in HW1, they (up to normalization) form a complete set of orthonormal basis for the Hilbert space $L_{\mathbb{C}^2}$ of linear operators. In the matrix representation, a generic state vector can be written as
\begin{align}
\dket{\psi}=\alpha\bpm\cos\theta e^{\imth\phi}\\ \sin\theta\epm,~~~\alpha\in\mathbb{C}
\end{align}
$\theta$ and $\phi$ are the two variables in the state vector that can be physically measured. One can carry out measurements for two observables to extract $\theta$ and $\phi$:
\begin{align}
\expval{\hat\sigma_z}=\cos(2\theta),~~~\expval{\hat\sigma_x}=\sin(2\theta)\cos\phi
\end{align}

\subsection{Appendix: Functions of operators/matrices}
\label{app:function of matrices}

A function $f(\hat M)$ of a square matrix $\hat M$ can be defined by the following Taylor expansion
\begin{align}
f(\hat M)=\sum_{n=0}^\infty\frac{\hat M^n}{n!}\frac{\text{d}^nf(x)}{\text{x}^n}|_{x=0}
\end{align}
if the Taylor expansion converges at $x=0$. In quantum mechanics we will mostly encounter square matrices that can be diagonalized, such as Hermitian and unitary matrices. 

For a diagonalizable matrix $\hat M$ with eigenkets $\{\dket{j}\}$ and eigenvalues $\{\lambda_j\}$, the matrix can be written as
\begin{align}
\hat M=\sum_j\lambda_j\dket{j}\dbra{j}
\end{align}
and a function of this diagonalizable matrix can be defined as
\begin{align}
f(\hat M)=\sum_jf(\lambda_j)\dket{j}\dbra{j}
\end{align}